\documentclass{my_memo}

\usepackage{booktabs}
\usepackage{hyperref}
\usepackage{xcolor}
\usepackage{listings}
\usepackage{longtable}
\sloppy

\lstset{
  language=Fortran,
  basicstyle=\ttfamily\small,
  keywordstyle=\color{blue!60!black},
  commentstyle=\color{green!40!black},
  stringstyle=\color{red!60!black},
  columns=fullflexible,
  breaklines=true,
  frame=single,
  framesep=4pt,
  showstringspaces=false,
  numbers=none
}
\lstdefinestyle{sh}{language=bash,keywordstyle=\color{black},basicstyle=\ttfamily\small}

\newcommand{\code}[1]{\texttt{#1}}
\newcommand{\um}{\ensuremath{\,\mu\mathrm{m}}}
\newcommand{\Cabs}{\ensuremath{C_{\rm abs}}}
\newcommand{\Cpol}{\ensuremath{C_{\rm pol}}}
\newcommand{\falign}{\ensuremath{f_{\rm align}}}

\title{Polarization in SEDust: Implementation Notes}
\author{Kwang-il Seon}
\date{2026 July 20}

\begin{document}
\maketitle

\begin{abstract}
\noindent
These are maintenance notes for the polarization capability added to SEDust:
where the polarized optics come from, how they are assembled, which
implementation choices were made and why, and how far the numbers have been
checked. The reader is assumed to be someone extending or repairing this code,
not someone calling it. The calling interface is documented in
\code{SEDust\_user\_manual.pdf}; the underlying physics is documented in
\code{aligned\_grain\_polarization.pdf}. Neither is repeated here.
\end{abstract}

\tableofcontents

%====================================================================
\section{Scope, and what lives elsewhere}

Three documents cover this work and they do not overlap
(Table~\ref{tab:whichdoc}). This one answers only the question \emph{how was
this built and why that way}. Where physics is needed to justify a coding
choice, the relevant section of the theory report is cited rather than
restated.

\begin{table}[h]
\centering
\begin{tabular}{@{}p{0.46\textwidth}p{0.46\textwidth}@{}}
\toprule
Question & Document \\
\midrule
Why polarization at all; alignment theory; the polarized transfer equation;
sign and index conventions; validation targets &
\code{aligned\_grain\_polarization.pdf} (\S1--\S5, \S9--\S11) \\[2pt]
How an RT code calls the library: signatures, units, what is already
integrated &
\code{SEDust\_user\_manual.pdf}, \S5 \\[2pt]
How the code is built, why each choice was made, what was measured &
this note \\
\bottomrule
\end{tabular}
\caption{Division of the polarization documentation.}
\label{tab:whichdoc}
\end{table}

\paragraph{Where this is going.} The polarization capability is not an end in
itself. The destination is the polarized radiative transfer of
Seon (2018): dust scattering plus dichroic extinction by aligned grains,
computed with grain optics derived from first principles, in the Monte Carlo
formalism of Peest et al.\ (2017) as extended to aligned spheroids by
Peest et al.\ (2023). Under that contract MoCafe supplies the alignment
assumptions (the \falign\ profile and the local magnetic-field geometry) and
SEDust returns every material quantity in matrix form: the $4\times4$
extinction matrix $K(\theta_i)$ with dichroism and birefringence, the phase
matrix $Z_{\rm al}(\theta_i;\theta_s,\phi)$ of the aligned population, the
random-orientation remainder, and the scattering cross sections. Dichroic
extinction, polarized emission, birefringence, and the orientation-resolved
cross-section table cover the extinction and emission side of that contract; the
scattering side --- the fixed-orientation Mueller matrix, the extinction matrix
built from forward amplitudes, and their library interface --- is now built and
verified, and is documented in \S\ref{sec:aligned}. What remains is MoCafe's
consumption: the frame rotations, direction sampling, peel-off, and the
$\exp(-K\tau)$ transfer step, none of which is a SEDust responsibility.

%====================================================================
\section{The data source}\label{sec:source}

Everything polarized derives from one file, the orientation-resolved
Draine \& Hensley (2021) astrodust spheroid table:
\begin{lstlisting}[style=sh]
data/dielectric/q_DH21Ad_P0.20_Fe0.00_1.400.dat.gz
\end{lstlisting}
(porosity $P=0.20$, iron fraction $f_{\rm Fe}=0$, axial ratio $b/a=1.400$).
Its layout is not self-describing, so it is recorded here
(Table~\ref{tab:filelayout}).

\begin{table}[h]
\centering
\begin{tabular}{@{}ll@{}}
\toprule
Property & Value \\
\midrule
header                & 12 lines, skipped \\
payload order         & \code{((Q(jw,jr,jori), jw=0,1128), jr=0,168), jori=1,3} \\
repeated              & three times: $Q_{\rm ext}$, then $Q_{\rm abs}$, then $Q_{\rm sca}$ \\
one record on disk    & the 169 sizes of one $(jw,\,jori)$ pair \\
records read          & $3~\text{quantities}\times3~\text{orientations}\times1129$ \\
total values          & $3\times3\times1129\times169 = 1{,}717{,}209$ \\
wavelength axis       & \code{data/dielectric/DH21\_wave}, 1129 nodes \\
size axis             & \code{data/dielectric/DH21\_aeff}, 169 nodes \\
axis file format      & two title lines, then the values free-format \\
\bottomrule
\end{tabular}
\caption{Layout of the orientation-resolved $Q$ table. The wavelength and size
axes are \emph{not} parsed out of the header; they come from the companion grid
files, which list the nodes the table was computed on.}
\label{tab:filelayout}
\end{table}

\paragraph{Orientation index.} The convention is the trap most likely to bite a
future reader, and it is worth stating in full (Table~\ref{tab:jori}); see also
\S3.6 of the theory report.

\begin{table}[h]
\centering
\begin{tabular}{@{}cll@{}}
\toprule
\code{jori} & geometry & meaning \\
\midrule
1 & $\mathbf{k}\parallel\hat{a}$        & propagation along the symmetry axis \\
2 & $\mathbf{k}\perp\hat{a}$, $\mathbf{E}\parallel\hat{a}$ & field along the symmetry axis \\
3 & $\mathbf{k}\perp\hat{a}$, $\mathbf{E}\perp\hat{a}$     & field across the symmetry axis \\
\bottomrule
\end{tabular}
\caption{Orientation index of the DH21 table. $\hat{a}$ is the spheroid
symmetry axis.}
\label{tab:jori}
\end{table}

%====================================================================
\section{Derived quantities}

For a grain whose symmetry axis lies in the plane of the sky and which is
perfectly aligned, the polarization cross section is the difference the grain
presents to the two linear polarizations, halved:
\begin{equation}
Q_{\rm pol} = \tfrac{1}{2}\left[\,Q(\code{jori}{=}3) - Q(\code{jori}{=}2)\,\right],
\label{eq:qpol}
\end{equation}
formed separately from $Q_{\rm ext}$ and from $Q_{\rm abs}$, giving
$Q_{\rm pol,ext}$ and $Q_{\rm pol,abs}$. The random-orientation average of the
same table is the trace,
\begin{equation}
Q_{\rm ran} = \tfrac{1}{3}\left[\,Q_1 + Q_2 + Q_3\,\right],
\label{eq:qran}
\end{equation}
which is the modified picket fence approximation (MPFA; \S3.2 of the theory
report). Cross sections follow from
\begin{equation}
C = Q\,\pi a_{\rm eff}^{2},
\end{equation}
with $a_{\rm eff}$ converted to cm. The alignment efficiency is the
Hensley \& Draine (2023) Table 1 fit,
\begin{equation}
\falign(a) = \frac{f_{\rm max}}{1 + (a_{\rm align}/a)^{\alpha}},
\qquad
a_{\rm align} = 0.0749\um,\;\;
\alpha = 1.80,\;\;
f_{\rm max} = 1.00 .
\label{eq:falign}
\end{equation}

%====================================================================
\section{Module map}

Table~\ref{tab:modules} lists what each file is responsible for. Nothing else
in the tree touches polarization.

\begin{table}[h]
\centering
\small
\begin{tabular}{@{}p{0.30\textwidth}p{0.63\textwidth}@{}}
\toprule
File / entity & Responsibility \\
\midrule
\code{sed/src/q\_table\_jori.f90} &
The reader. Decompresses, parses the three quantity blocks, validates
(monotone axes, finite values, no trailing data), forms
Eqs.~(\ref{eq:qpol})--(\ref{eq:qran}), and supplies \code{falign\_hd23(a)},
Eq.~(\ref{eq:falign}). Publishes \code{qpol\_ext}, \code{qpol\_abs} and
\code{qran\_*}; the orientation-resolved \code{qext\_j}/\code{qabs\_j}/%
\code{qsca\_j} are \code{private} and are deallocated once the combinations
have been formed. \\[2pt]
\code{sed\_astrodust.f90}: \code{build\_Cpol} &
Interpolates $Q_{\rm pol}$ from the 169-node table axis onto the
size-distribution grid in $\log a$, exactly as \Cabs\ and $C_{\rm sca}$ are,
multiplies by $\pi a_{\rm eff}^2$, and fills \code{falign\_ad}. Checks that the
polarized and $Q$ wavelength grids agree node for node. \\[2pt]
\code{sed\_astrodust.f90}: \code{Cpol}, \code{Cpol\_ext}, \code{falign\_ad} &
Module-level arrays, \code{(NLAM,NA)} and \code{(NA)}. \code{Cpol} is the
\emph{absorption} one and drives polarized emission; \code{Cpol\_ext} is its
extinction twin and drives dichroic extinction. \\[2pt]
\code{dust\_model\_mod.f90}: \code{grain\_pop\_t} &
Carries \code{Cpol(:,:)} and \code{falign(:)}. Both unallocated is how the
solver recognizes an unpolarized population. Also carries the extinction-side
\code{Cpol\_ext(:,:)} and \code{gsca(:,:)}, read only by
\code{dust\_extinction}. \\[2pt]
\code{sed\_astrodust.f90}: \code{set\_pop} &
Optional \code{Cpol\_in}, \code{falign\_in}. Both must be supplied together;
supplying neither leaves the population unpolarized. The extinction-side
\code{Csca\_in}, \code{Cpol\_ext\_in} and \code{gsca\_in} are optional and
independent; a population that omits them contributes zero to those terms of
the size integral. \\[2pt]
\code{sed\_astrodust.f90}: \code{sed\_grain\_loop} &
Optional \code{Cpol\_pop}, \code{falign\_pop}, \code{Jpol}, again all three or
none (\code{do\_pol}). Accumulates the polarized channel alongside \code{Jout}
in every solver branch (\S\ref{sec:sites}). \\[2pt]
\code{sed\_astrodust.f90}: \code{dust\_emission} &
Optional \code{lamI\_pol}, last argument. Sums \code{Jpol} over the populations
that carry polarized optics and applies the same unit convention as
\code{lamI\_total}. \\[2pt]
\code{sed\_astrodust.f90}: \code{dust\_extinction} &
The extinction accessor, per H atom on \code{m\%lam}. It integrates
$C_{\rm pol,ext}\,\falign$ (and the birefringent twin) over the size
distribution of every population --- these are the same numbers column 8 of the
table carries, without the file. \Cabs, $C_{\rm sca}$ and
$\langle\cos\theta\rangle$ it later stopped integrating: it reads them from the
precomputed \code{kext} table instead, while the polarized outputs stay
computed because their \falign\ weight is runtime state
(\S\ref{sec:whyaccessor}). The first-principles integral of all six is
\code{size\_integrated\_extinction}. \\[2pt]
\code{sed/src/calc\_polext.f90} &
Standalone check: computes the polarized extinction directly from
\code{qpol\_ext} and the release size distribution, and compares against the
published \code{polarized\_extinction.dat}. Independent of \code{sed\_init}. \\[2pt]
\code{sed/src/calc\_kext.f90} &
The extinction-table driver, one executable for every dust model
(\code{./calc\_kext.x astrodust\,|\,dl07\,|\,zubko\,|\,from\_files}). In its
\code{astrodust} mode it writes column 8, \code{C\_polext/H}, of
\code{data/kext\_astrodust\_MW.dat} as
$\sum_a {\rm d}n_{\rm Ad}(a)\,C_{\rm pol,ext}(\lambda,a)\,\falign(a)$. Only that
mode writes the column, because only the astrodust builder fills
\code{Cpol\_ext}. \\
\bottomrule
\end{tabular}
\caption{Responsibility of each file touched by the polarization update.}
\label{tab:modules}
\end{table}

%====================================================================
\section{Implementation decisions}

This is the part worth reading again in six months.

\subsection{The accumulation sites, and the two that hide}\label{sec:sites}

\code{sed\_grain\_loop} branches on \code{stoch\_method} into four cases, and
the polarized channel has to be accumulated wherever the total is. There are
eight such statements (Table~\ref{tab:sites}). Missing any one of them produces
a polarized spectrum that is silently too low rather than an error.

\begin{table}[h]
\centering
\small
\begin{tabular}{@{}llll@{}}
\toprule
branch & context & accumulates into & pattern \\
\midrule
\code{'draine'}    & equilibrium grain            & \code{Jpol\_local} & thread-local \\
\code{'draine'}    & stochastic, sum over $P(T)$  & \code{Jpol\_local} & thread-local \\
\code{'heuristic'} & equilibrium grain            & \code{Jpol}        & \textbf{direct} \\
\code{'heuristic'} & stochastic, sum over $P(T)$  & \code{Jpol}        & \textbf{direct} \\
\code{'qm'}        & QM solve, rescaled           & \code{Jpol\_local} & thread-local \\
\code{'qm'}        & GD fallback after QM failure & \code{Jpol\_local} & thread-local \\
\code{'qm'}        & equilibrium grain            & \code{Jpol\_local} & thread-local \\
\code{'equil'}     & every grain                  & \code{Jpol}        & direct \\
\bottomrule
\end{tabular}
\caption{Where the polarized emission is accumulated. The two boldface rows
are the ones a search for the \code{Jout\_local} pattern does not find.}
\label{tab:sites}
\end{table}

\paragraph{The lesson.} The \code{'draine'} and \code{'qm'} branches are
OpenMP-parallel and accumulate into a thread-local \code{Jout\_local} that is
reduced at the end. The \code{'heuristic'} branch is serial and accumulates
into \code{Jout} directly. A survey of the accumulation sites conducted by
searching for \code{Jout\_local} therefore returns only the parallel branches,
and an initial pass over this code identified four sites on exactly that basis.
Since \code{'heuristic'} is the \emph{default} solver, that omission would have
made the polarized emission come back identically zero in the default
configuration while every other solver looked correct. Any future change that
adds a term to the emission sum must be checked against
Table~\ref{tab:sites} branch by branch, not by pattern search.

\subsection{Why the \code{'qm'} rescaling is exact}

\code{qm\_solve\_grain} does not return a temperature distribution; it returns
the emitted spectrum with \Cabs\ already folded in. The polarized counterpart
is therefore obtained by rescaling:
\begin{lstlisting}
if (Cabs_pop(ii, ir) > 0.0_wp) then
   Jpol_local(ii) = Jpol_local(ii) + &
      wpol * emission_qm(ii) * (Cpol_pop(ii, ir) / Cabs_pop(ii, ir)) / &
      (4.0_wp * PI * lam(ii) * 1.0e-3_wp)
end if
\end{lstlisting}
This is not an approximation. The temperature distribution $P(T)$ is set by the
\emph{heating}, which is governed by \Cabs\ and is unchanged by the question of
which polarization is being read out. Only the emission readout changes, so
\begin{equation}
C_{\rm pol}\!\int\! P(T)\,B_\lambda(T)\,{\rm d}T
=\frac{C_{\rm pol}}{C_{\rm abs}}\;
 C_{\rm abs}\!\int\! P(T)\,B_\lambda(T)\,{\rm d}T
=\frac{C_{\rm pol}}{C_{\rm abs}}\;\code{emission\_qm}
\end{equation}
holds identically, wavelength by wavelength. Wavelengths with zero \Cabs\ emit
nothing and are skipped rather than divided by zero.

\subsection{The $P(T)$ solver was not touched}

No change was made to \code{calc\_P}, \code{narrow\_iterative}, or
\code{qm\_solve\_grain}. This is deliberate and it is what makes stochastic
heating automatically correct for polarization: at every site the same
temperature weight multiplies both accumulators, one with
${\rm d}n\,\Cabs$ and the other with ${\rm d}n\,\falign\,\Cpol$. The polarized
spectrum is thus the same average over the same $P(T)$, and no separate
stochastic treatment of the polarized channel is needed (see \S5.7 of the
theory report for why this is the physically right structure).

\subsection{Mixed orientation-average conventions}\label{sec:mixed}

This is the one open physical wart, and it is documented rather than hidden.

\Cabs\ in SEDust comes from our own T-matrix run, whose random-orientation
average is computed exactly. \Cpol\ comes from the release table, whose
available average is the $1/3$ trace, Eq.~(\ref{eq:qran}), i.e.\ the MPFA. The
two are different approximations to the same quantity, and the code mixes them:
\code{build\_Cpol} adds the polarized channel and deliberately does \emph{not}
recompute \Cabs.

Table~\ref{tab:mixed} gives the measured disagreement,
$|Q_{\rm abs}^{\rm trace}/Q_{\rm abs}^{\rm exact}-1|$, on the shared
$(\lambda, a_{\rm eff})$ grid.

\begin{table}[h]
\centering
\begin{tabular}{@{}lrr@{}}
\toprule
band & median & maximum \\
\midrule
far infrared ($\lambda > 30\um$)     & 0.022\% & 3.46\% \\
\quad restricted to $a \le 0.5\um$   &         & 1.98\% \\
mid infrared                         & 0.033\% &        \\
near infrared                        & 0.038\% &        \\
optical                              & 0.080\% &        \\
ultraviolet                          & 0.21\%  & 9.53\% \\
\bottomrule
\end{tabular}
\caption{Disagreement between the trace average and the exact
orientation average of $Q_{\rm abs}$.}
\label{tab:mixed}
\end{table}

\paragraph{Reading the far-infrared maximum.} The 3.46\% worst case occurs at
$a = 4.2\um$, a size carrying ${\rm d}n/n_{\rm H} \sim 2.8\times10^{-36}$, so
it is numerically irrelevant. Restricting to $a \le 0.5\um$, which covers
99.999\% of the geometric cross section, the worst case drops to 1.98\%.
At $a \ll \lambda$ both averages approach the same Rayleigh limit, which is why
the far infrared is well behaved.

\paragraph{Verdict: deliberately deferred in this release.} The mixture is
left in place for v1.20, as a decision rather than an oversight. The reasoning
is recorded here so that it is not rediscovered later. The mixture is harmless
for polarized \emph{emission},
which is a far-infrared phenomenon. It is \textbf{not} defensible for
polarized \emph{extinction} in the ultraviolet or the optical, and column 8
spans the whole wavelength grid. The \code{Cpol\_ext} output of
\code{dust\_extinction} is the same quantity by the same route and carries the
same caveat unchanged. Anyone using either shortward of the near
infrared must resolve this first. The clean resolution is to recompute \Cabs\
on the MPFA convention so that both channels share one average. That changes
the existing SEDs, so it breaks their byte-identity with the current reference
outputs and has to be done as a separate, deliberately validated step rather
than folded into this one. The size of the change is not an argument for
leaving it: the correctness argument is that two channels of one calculation
should not use two different orientation averages.

What was decided for this release is only \emph{when}, not \emph{whether}.
The planned use of these optics --- redoing Seon (2018) in the $V$ band, with
one or two further optical bands possible --- lies entirely in the range where
the disagreement is $0.080\%$ or below, while unifying the convention would
change \Cabs\ and break the byte-identical regression baseline that the rest
of v1.20 was validated against. Spending that baseline buys nothing inside the
present scope. It has to be spent before any ultraviolet polarized extinction
is quoted, and that is the trigger for doing the work.

\subsection{\falign\ is the fit, not the release column}

\code{data/release/size\_distribution.dat} also ships an alignment column. It
agrees with Eq.~(\ref{eq:falign}) to a median ratio of 1.0000 but differs by up
to 0.32\% at the small end. The analytic form is the canonical one in SEDust
because it is the published fit, because it is defined on any radius grid
rather than only the tabulated one, and because it keeps the alignment model in
a single place. Do not mix the two sources: a quantity computed half from one
and half from the other is neither.

\subsection{Names and surface area}

\paragraph{\code{Cpol} was not renamed to \code{Cpol\_abs}.} The name appears at
roughly 40 sites and is part of the public surface
(\code{grain\_pop\_t\%Cpol}, the optional \code{set\_pop(..., Cpol\_in=...)}).
A rename would touch every one of them for no functional gain. The declaration
carries an explicit note instead: \code{Cpol} is the absorption one,
\code{Cpol\_ext} its extinction twin.

\paragraph{The polarized extinction was exported as a table column first.} A
dedicated \code{dust\_cpol\_ext} accessor was considered at the time and
rejected: \code{dust\_lib} was then an emission-facing API with no extinction
surface at all, so such an accessor would have been the module's only
extinction entry point, a worse inconsistency than the inconvenience it
removed. The condition attached to that decision was that \code{Cpol\_ext}
should join a full extinction API if one ever appeared. It has;
see \S\ref{sec:whyaccessor}.

\subsection{Why the accessor exists now}\label{sec:whyaccessor}

The table-only export was right for a library that emitted and did not
extinguish. Rebuilding a radiative transfer code on top of SEDust changed the
premise: the host would have obtained its emission from a call and its opacity
from a parsed text file, and that asymmetry is the whole problem. The file is
also written on whatever wavelength grid the driver used, so a host on any
other grid has to interpolate optics that the model could have evaluated
exactly.

\code{dust\_extinction} closes it. The model object already held everything
required, so the accessor adds no physics and no new data source; it performs
the same size integral the driver performs, over the populations already
built. Two fields were added to \code{grain\_pop\_t} to make that possible:
\code{Cpol\_ext}, previously a module array only, and \code{gsca}, the
scattering asymmetry, which the emission path never needed and which
\code{sed\_init} now interpolates onto the size grid beside \Cabs\ and
$C_{\rm sca}$, from the same $Q$ table and by the same routine the driver used.
Agreement with the table is at the precision the file is written with (a median
relative deviation of $\sim7\times10^{-9}$ in each cross section), which is the
expected result for one integral evaluated twice.

The table remains, unchanged byte for byte, and remains the right route for a
tool that wants the optics without linking the library.

\paragraph{Later revision: the scalar optics went back to the table.} The
argument above is about the \emph{interface}, and it stands: a host must get
its opacity from a call on the same model object that gives it emission, on its
own wavelength grid. What changed afterwards is where that call gets its
numbers. The scalar transport optics of a dust model do not depend on the
transport, so recomputing the size integral inside an RT run buys nothing that
reading the recorded integral does not. \code{dust\_extinction} therefore now
serves \Cabs, $C_{\rm sca}$, $C_{\rm ext}$ and $\langle\cos\theta\rangle$ from a
\code{data/kext\_*.dat} table, attached to the model by the builder (optional
\code{kext\_path}; the defaults are the EUV products, the widest grid each model
has) and interpolated onto \code{m\%lam}, exactly at the nodes the two grids
share. The host's call site is unchanged, and so is the requirement that the
file be opened while the working directory is still \code{sed/} --- which is why
the read happens in \code{build\_*} and not in \code{dust\_extinction}.

The polarized outputs were deliberately \emph{not} moved. \code{Cpol\_ext} and
\code{Cbir\_ext} carry the \falign\ weight, which a host resets cell by cell
through \code{dust\_set\_alignment}; a column fixed at build time could not
follow it. So the routine is now mixed on purpose: scalar optics from a table,
polarized optics from the model. The size integral itself is unchanged and is
still reachable, for all six outputs, as
\code{size\_integrated\_extinction} --- which is what \code{calc\_kext.x} calls
to write the tables, and what the in-tree checks call when they must compare one
build of a model against another.

\subsection{Compression handling, and a scratch file that leaked}

The table ships compressed and Fortran cannot read a deflate stream. Adding a
zlib binding was rejected as too heavy for one file. The reader instead shells
out once,
\begin{lstlisting}
call execute_command_line('gzip -dc "'//trim(gz_file)//'" > "'// &
                          trim(out_file)//'"', exitstat=estat, cmdstat=cstat)
\end{lstlisting}
decompresses to a scratch copy, reads it, and deletes it. A fresh clone works with
no manual setup step, and neither a 17\,MB untracked sibling in
\code{data/dielectric} nor a new build dependency is created. A path not ending in
\code{.gz} is opened directly.

\paragraph{The scratch target: \code{TMPDIR}, with the process id.} The scratch
copy is written under \code{TMPDIR} (else \code{/tmp}), named with the process id:
\code{q\_jori\_<pid>.dat} for the polarized $Q$ reader (\code{q\_table\_jori}) and
\code{scatmat\_aligned\_<pid>.dat} for the aligned scattering reader
(\code{scatmat\_aligned\_mod}); both loaders share the scheme. Naming by process id
means concurrent MPI ranks never write the same scratch file, and placing it in
\code{TMPDIR} rather than the launch directory means a read-only launch directory
is never written to.

\paragraph{The bug this had.} An earlier version wrote a fixed
\code{q\_jori\_scratch.dat} in the \emph{working} directory. Shell redirection
creates the target file \emph{before} \code{gzip} runs, so when decompression
failed the early return left a zero-length \code{q\_jori\_scratch.dat} behind
there. The failure path now calls \code{discard\_scratch\_copy} before bailing
out, like every other error exit --- and, since the scratch file moved to
\code{TMPDIR/q\_jori\_<pid>.dat}, a leaked remnant no longer lands in the user's
directory at all. The general shape of the bug, an error path that skips the
removal the success path performs, is worth watching for in any future edit to
this reader: it has eight separate early returns.

%====================================================================
\section{Verification}

\subsection{Polarized extinction against the release}

\begin{table}[h]
\centering
\begin{tabular}{@{}ll@{}}
\toprule
Check & Result \\
\midrule
Column 8 vs.\ released \code{polarized\_extinction.dat} &
median 0.0296\%, max 0.9422\%, 992 points \\
Column 8 vs.\ \code{calc\_polext.x} &
agree to $7\times10^{-8}$ relative \\
Columns 1--7 of \code{kext\_astrodust\_MW.dat} &
byte-identical to the pre-change file \\
\code{astrodust\_irem\_ours\_S1/S2/PAH.dat} &
byte-identical \\
\code{dl07\_sed\_ours\_mw31\_60.dat} &
byte-identical \\
Build & zero warnings \\
\bottomrule
\end{tabular}
\caption{Verification of the polarized extinction path and regression of
everything else.}
\label{tab:verif}
\end{table}

The comparison against the release is restricted to $\lambda \ge 0.11\um$
because the reference file turns negative below that. The 0.9422\% maximum
deviation is not a systematic offset: it sits at $\lambda = 0.112\um$, right at
the sign reversal of the polarized extinction, where the quantity passes
through zero and a relative deviation is not a meaningful measure. Away from
that point the agreement is at the level of the median.

The $7\times10^{-8}$ figure is the precision the ASCII output format carries,
so column 8 and \code{calc\_polext.x} agree as closely as the files allow.
They are two independent routes to the same quantity: \code{calc\_polext.x}
works straight from \code{qpol\_ext} and the release size distribution, while
column 8 goes through \code{sed\_init}, the $\log a$ interpolation onto the SED
size grid, and \code{build\_Cpol}.

\subsection{Polarized emission}

Intrinsic polarization fraction \code{lamI\_pol/lamI\_total} under a Mathis
ISRF at $U = 1.585$:

\begin{table}[h]
\centering
\begin{tabular}{@{}lr@{}}
\toprule
$\lambda$ & fraction \\
\midrule
12\um                       & 0.00\% \\
100\um                      & 13.67\% \\
154\um                      & 17.15\% \\
350\um                      & 18.79\% \\
850\um                      & 19.15\% \\
median over 300--3000\um    & 19.18\% \\
\bottomrule
\end{tabular}
\caption{Intrinsic polarization fraction of the emitted spectrum, before any
geometric factor.}
\label{tab:polfrac}
\end{table}

The mid-infrared zero is correct rather than a failure: there the emission is
dominated by stochastically heated PAHs, which HD23 take to be unaligned.

\paragraph{19.15\% versus Planck's 22\% is a property of the model, not a bug.}
Planck measures $p_{\rm max} = 22.0\,(+3.5,-1.4)\%$ at 353\,GHz. The 19.15\%
here is the \emph{maximum} attainable value, obtained before any geometric
factor, and every real line of sight reduces it, so the model as implemented
cannot reach the observed maximum. This traces to the dust model itself: the
axial ratio $b/a = 1.4$ is the minimum non-sphericity DH21b found the starlight
polarization efficiency to require, and it was not tuned to maximize the
submillimeter polarization fraction. Do not close the gap by adjusting
\falign\ or the axial ratio without recognizing that as a change to the dust
model, requiring its own revalidation against the extinction and emission
constraints HD23 fit.

%====================================================================
\section{First-principles orientation-resolved optics}\label{sec:firstprinciples}

The table of \S\ref{sec:source} is Draine \& Hensley's precomputed one. SEDust
can now regenerate the same orientation-resolved cross sections from the
astrodust dielectric function with its own T-matrix engine, so the polarized
optics no longer have to be taken on trust from the release file. The
random-orientation optics were already SEDust's own (the \code{q\_astrodust}
table read by \code{q\_table}); this closes the remaining gap, the
orientation-resolved part read by \code{q\_table\_jori}.

Recall the three orientations of the incident wave relative to the spheroid
symmetry axis $\hat a$: $\mathrm{jori}=1$ is $\mathbf{k}\parallel\hat a$;
$\mathrm{jori}=2$ is $\mathbf{k}\perp\hat a$ with $\mathbf{E}\parallel\hat a$;
$\mathrm{jori}=3$ is $\mathbf{k}\perp\hat a$ with $\mathbf{E}\perp\hat a$. The
polarized and mean combinations are
$Q_{\rm pol}=\tfrac12(Q_3-Q_2)$ and $Q_{\rm ran}=\tfrac13(Q_1+Q_2+Q_3)$.

The optics are reached two ways. The \emph{table} is loaded once and
interpolated: this is the production path, fast and safe to call from many
threads. The \emph{direct} routine \code{oriented\_cross\_sections} computes one
node from first principles; it is the shared core the table generator loops
over, and it also serves an arbitrary $(a,\lambda)$ or a shape or porosity for
which no table exists. It must not be called from inside a loop over
radiative-transfer cells: each call is a full T-matrix solve plus an angular
integral, and it works through an initialized T-matrix workspace of 80\,MiB.
Thread safety is no longer the objection --- the library holds every working
array in a caller-owned \code{tmatrix\_workspace\_t}, with no \code{COMMON}, no
\code{EQUIVALENCE} and no mutable \code{SAVE} (\S\ref{sec:aligned:gen}) --- the
cost is.

\subsection{Rayleigh regime, $x<0.1$}

In the electrostatic limit the oblate spheroid has a diagonal polarizability
$\boldsymbol\alpha=\mathrm{diag}(\alpha_a,\alpha_b,\alpha_b)$, with $\alpha_a$
for $\mathbf{E}$ along the symmetry axis and $\alpha_b$ for $\mathbf{E}$
transverse, set by the depolarization factors $L_a,L_b$ of the spheroid through
$\alpha_i\propto(\varepsilon-1)/[\,1+(\varepsilon-1)L_i\,]$. Absorption and
scattering follow as $Q_{\rm abs}\propto\mathrm{Im}\,\alpha$ and
$Q_{\rm sca}\propto|\alpha|^2$. The three orientations sample exactly these two
polarizabilities: $\mathrm{jori}=1$ and $3$ present the transverse $\alpha_b$
(the field is transverse to $\hat a$ in both), $\mathrm{jori}=2$ presents
$\alpha_a$. So the orientation-resolved cross sections are the two polarization
components the random average $(\alpha_a+2\alpha_b)/3$ is already built from;
extracting them before the average is analytic and exact, and $Q(\mathrm{jori}=1)$
equals $Q(\mathrm{jori}=3)$ identically. \code{rayleigh\_limit} returns them
through optional arguments, leaving its averaged outputs untouched.

\subsection{T-matrix regime, $0.1\le x\le 50$}

The T-matrix $\mathbf{T}$ is a property of the particle alone --- size, shape,
refractive index --- independent of orientation and of the incidence and
scattering directions. SEDust solves it once for each wavelength and radius with
\code{tmatrix\_eval\_scattering\_matrix}, the library entry point that descends
from Mishchenko's random-orientation solver; the converged $\mathbf{T}$ stays in
the caller's workspace and the returned \code{tmatrix\_scatmat\_t} identifies it.
Two orientation-resolved quantities are then read off the same $\mathbf{T}$.

\paragraph{Extinction, from the optical theorem.} Mishchenko's \code{AMPL}
evaluates the fixed-orientation amplitude matrix
$\mathbf{S}(\mathbf{k}\!\to\!\mathbf{k}')$ from the stored $\mathbf{T}$ for a
particle placed at chosen Euler angles. Taking the scattering direction equal to
the incidence direction (forward) and applying the optical theorem,
$C_{\rm ext}=(4\pi/k)\,\mathrm{Im}\,S_{\rm fwd}$ for the incident polarization,
gives $C_{\rm ext}$ for each orientation: the symmetry axis is placed along
$\mathbf{k}$ ($\mathrm{jori}=1$) or across it ($\mathrm{jori}=2,3$) and the
co-polar forward amplitude is read ($S_{\theta\theta}$ for $\mathbf{E}$ along
$\hat\theta$, $S_{\phi\phi}$ for $\hat\phi$).

\paragraph{Scattering, from the phase-matrix integral.} The same \code{AMPL},
swept over all scattering directions at fixed incidence, gives the differential
scattering cross section $|\mathbf{S}|^2$; the scattering cross section is its
integral over the sphere,
$C_{\rm sca}=\int(|S_{\theta\leftarrow p}|^2+|S_{\phi\leftarrow p}|^2)\,
d\Omega$ for incident polarization $p$, and absorption follows by conservation,
$C_{\rm abs}=C_{\rm ext}-C_{\rm sca}$. The fixed-orientation scattered intensity
is a trigonometric polynomial of degree ${\sim}2N_{\max}$ in $\cos\theta$ and in
$\phi$, so Gauss--Legendre nodes in $\cos\theta$ ($N_{\max}+2$ of them) and a
uniform azimuth grid ($2N_{\max}+2$ points) integrate it to machine precision.

\paragraph{An engine subtlety, and how it was retired.} In Mishchenko's f77
sources the random-orientation expansion \code{GSP} reuses the \code{/TMAT/}
storage as scratch through an \code{EQUIVALENCE} and so destroys the T-matrix it
reads. The amplitude evaluation needs that same T-matrix afterward, so the f77
driver kept an intact copy in a second common block and restored \code{/TMAT/}
after \code{GSP}. Without that restore the oriented cross sections came out as
garbage (negative $C_{\rm ext}$), which corrected a premise stated in the
planning note, that the converged T-matrix simply remains available after the
solve. The library removes the overlay rather than repairing it: the workspace
carries the T-matrix (\code{full\_tstore}) and the expansion scratch
(\code{full\_gsp\_*}) as separate components, the expansion routine
(\code{GSP\_FULL\_DIRECT}) only reads the former, and it therefore survives the
solve intact with nothing to restore.
The planning note's premise holds again, now by construction; the episode is
kept on the record because the failure it caused was silent.

\paragraph{Reusing a workspace is checked, not trusted.} A workspace holds one
T-matrix at a time, so an amplitude evaluation made against a \code{scatmat}
older than the T-matrix now in the workspace would silently answer for the wrong
particle. The workspace carries a generation counter \code{tm\_tag} that advances
on every evaluation, successful or failed, and
\code{tmatrix\_eval\_scattering\_matrix} copies it into
\code{scatmat\%state\_tag}; \code{AMPL} compares the two and refuses with
\code{IERR = 4} when they differ. The remaining codes are \code{0} success,
\code{1} an angular argument outside its range, \code{2} no converged T-matrix in
the workspace, \code{3} a truncation order outside the stored block. Nothing here
is a convenience: each code replaces a case that used to be undefined behavior or
a \code{STOP} inside a library.

\subsection{Geometric-optics regime, $x>50$}

Here the T-matrix does not converge and an asymptotic model is used. Extinction
is the extinction paradox, $C_{\rm ext}\to 2A_{\rm proj}$, twice the geometric
shadow. For an oblate spheroid of equatorial semi-axis $a_s$ and symmetry
semi-axis $c$ ($a_s/c=b/a$), the shadow is $\pi a_s^2$ when $\mathbf{k}$ is along
the symmetry axis ($\mathrm{jori}=1$) and $\pi a_s c$ when it is across
($\mathrm{jori}=2,3$), giving $Q_{\rm ext}(1)=2(b/a)^{2/3}$ and
$Q_{\rm ext}(2)=Q_{\rm ext}(3)=2(b/a)^{-1/3}$. The shadow depends on the
direction of $\mathbf{k}$ relative to the axis but not on the incident
polarization, so extinction carries no polarization at this order.

The polarization in this regime is carried by absorption. In the opaque limit,
where a refracted ray is fully absorbed --- valid when
$\mathrm{Im}(m)\,x\gg1$ --- the absorptivity of an illuminated surface element is
one minus its Fresnel power reflectance, and that reflectance is
polarization-dependent ($R_s\neq R_p$ away from normal incidence). Integrating
over the illuminated half of the spheroid,
\[
  C_{{\rm abs},p}=\int_{\rm illum}
    \bigl[\,1-\bigl(|\hat e\!\cdot\!\hat s|^2 R_s(\theta_i)
                   +|\hat e\!\cdot\!\hat p|^2 R_p(\theta_i)\bigr)\bigr]
    \,dA_{\rm proj},
\]
where $\theta_i$ is the local angle of incidence and $\hat s,\hat p$ the local
senkrecht and parallel directions. The two grain-frame polarizations
($\hat e=\hat a$ versus $\hat e\perp\hat a$) project differently onto the local
$\hat s,\hat p$ basis as they sweep the curved surface, and that is the whole
source of the dichroism. A purely geometric ray treatment, with no Fresnel
coefficients, would give no polarization at all; the surface reflectance is what
makes the regime dichroic, and it is needed to match the release, which carries
a real absorption dichroism there --- about $3\%$ of the mean, and
\emph{growing} with size parameter, the wave-optics signature a ray limit cannot
reproduce. This is a deliberate approximation, documented at the code site with
its opaque-limit domain of validity; it would matter for a size distribution
with large-grain weight, which the astrodust distribution does not have.

\subsection{Dispatch, generation, and drop-in}

\code{oriented\_cross\_sections} selects the regime by $x$ and returns the
oriented $(Q_{\rm ext},Q_{\rm abs},Q_{\rm sca})$ for one node, with the same
branch boundaries and fallbacks as the random-orientation driver. \code{run\_q\_jori.x}
loops it over the DH21 grid and writes the table in the exact stream the release
reader expects (\S\ref{sec:source}), so the result is a drop-in:
\code{sed\_init} and \code{build\_astrodust} already take an optional
\code{qpol\_path} that defaults to the release file, and pointing it at the
regenerated table switches the polarized optics to SEDust's own with no code
change. A full sweep is dominated by the T-matrix nodes (of order 16~hours on
one core), so the driver has a wavelength-window \code{range} mode for
process-level parallelism and a \code{merge} mode that reassembles the windows
--- a window is not a contiguous slice of the stream, because the wavelength
index runs innermost.

\subsection{How far it has been checked}\label{sec:fpcheck}

Each regime was compared node by node against the release table
(Table~\ref{tab:fpcheck}). Across the Rayleigh and T-matrix regimes --- every
wavelength and radius that carries astrodust weight --- the reconstruction
matches the release to a median of a few parts in $10^4$, in both the mean and
the polarized cross sections. The geometric-optics region agrees only coarsely,
as an asymptotic model should: the mean absorption to about ten percent and the
dichroism to about three-quarters of the release value, both with the correct
sign and size-parameter trend. It joins the T-matrix continuously across $x=50$
in sign and ordering, and it carries negligible weight in the size integral.

End to end, feeding the regenerated table through \code{calc\_polext} reproduces
the released polarized extinction to a median of $0.06\%$ over
$\lambda\ge0.11\um$, against $0.03\%$ for the release table itself; the two agree
with each other to a median of $0.02\%$ at the level of the observable. So the
few-parts-in-$10^4$ reconstruction error at the level of a single cross section
survives the size integral as a few parts in $10^4$ in the polarized extinction,
computed entirely from the astrodust dielectric function with no recourse to the
release optics.

\begin{table}[h]
\centering
\small
\begin{tabular}{llll}
\toprule
Regime & $x$ & Quantity & median $|{\rm ours}/{\rm HD23}-1|$ \\
\midrule
Rayleigh  & $<0.1$      & $Q_{\rm ran}$ (abs) & $2.3\times10^{-4}$ \\
          &             & $Q_{\rm pol}$ (abs) & $3.7\times10^{-4}$ \\
T-matrix  & $0.1$--$50$ & $Q_{\rm ran}$ (ext) & $2.1\times10^{-4}$ \\
          &             & $Q_{\rm pol}$ (ext) & $8.4\times10^{-4}$ \\
          &             & $Q_{\rm ran}$ (abs) & $1.9\times10^{-4}$ \\
          &             & $Q_{\rm pol}$ (abs) & $7.6\times10^{-4}$ \\
GO        & $>50$       & $Q_{\rm ran}$ (abs) & $9.3\times10^{-2}$ \\
          &             & $Q_{\rm pol}$ (abs) & $3.3\times10^{-1}$ \\
\bottomrule
\end{tabular}
\caption{Node-by-node reconstruction against the HD23 release table. The
geometric-optics dichroism is at $\sim0.75$ of the release with matching sign
and trend; its size-parameter region carries no astrodust weight.}
\label{tab:fpcheck}
\end{table}

\subsection{Computing named wavelengths}\label{sec:namedlam}

Polarized transfer is run at a handful of wavelengths, not on a whole axis. The
\code{range} mode cannot serve that: its \code{JW1}/\code{JW2} are \emph{indices}
into a fixed wavelength file, so it reaches only wavelengths that file already
carries. Three modes take the physical wavelength instead:
\begin{lstlisting}[style=sh]
./run_q_jori.x lam L1 [L2 ...] [ja=JA1:JA2] [tag=NAME]   # wavelengths [um]
./run_q_jori.x lamfile PATH    [ja=JA1:JA2] [tag=NAME]   # one per line, '#' comments
./run_q_jori.x lammerge STEM FILE [FILE ...]             # reassemble ja= windows
\end{lstlisting}
Each writes a pair, \code{q\_astrodust\_jori\_P0.20\_Fe0.00\_1.400.TAG.dat} and
its wavelength axis \code{...TAG.wave}; the default \code{TAG} is \code{lamL} for
one wavelength and \code{lamLMIN\_LMAX\_nN} for several. That pair is what
\code{sed\_init} / \code{build\_astrodust} take through \code{qpol\_euv\_path} and
\code{qpol\_euv\_wave\_path}. The reader interpolates in $\log\lambda$, so at
least \emph{two} wavelengths are needed to fill a band; a one-wavelength file is
a valid product to read directly but cannot be interpolated.

\paragraph{Processes, not threads.} \code{ja=JA1:JA2} restricts one invocation to
radii \code{JA1..JA2} of the 169-point size axis so that several \emph{processes}
share one wavelength, and \code{lammerge} puts the column slices back side by
side. This was once forced by the engine: the f77 solver handed the converged
T-matrix to \code{AMPL} through \code{COMMON /TMAT/} and kept its working storage
in further \code{COMMON} blocks, two of them blank, so the core was not
re-entrant. That constraint is gone --- the library owns no shared state and
distinct workspaces may be evaluated concurrently
(\S\ref{sec:aligned:gen}) --- but the driver keeps the process windows, which are
already measured and cost nothing to run. Separate processes each get their own
copy, and a merged file has been checked byte for
byte against a single-process one. Measured: three wavelengths (0.0124, 0.0602 and
0.0912\um) over 57 processes of three radii each took 9\,m\,22\,s of wall time for
34\,m\,48\,s of processor time. The speedup is 3.7 rather than 57 because the cost
is not spread over the radii --- a T-matrix solve grows like $N_{\max}^4$ with
$N_{\max}\sim x$, so the few radii just below the $x = 60$ ceiling carry nearly
all of the work. One radius per invocation (169 processes) lets the operating
system interleave them and brings the wall time down to the single most expensive
node.

\paragraph{\code{euv} is an axis switch, not a wavelength.} The leading
\code{euv} keyword selects a wavelength \emph{axis file}
(\code{data/dielectric/DH21\_wave\_euv}), so it cannot be combined with
\code{lam}, \code{lamfile} or \code{lammerge}, which carry their own axis. The
driver rejects the combination explicitly.

\subsection{The polarized band below the Lyman limit}\label{sec:euvpol}

\Cpol, $C_{\rm pol,ext}$ and $C_{\rm bir,ext}$ are read from an
orientation-resolved table on the full DH21 wavelength axis, all 1129 nodes from
$0.0912$ to $39810\um$: by default the release file
\code{data/dielectric/q\_DH21Ad\_P0.20\_Fe0.00\_1.400.dat.gz}, or, through
\code{qpol\_path}, the regenerated
\code{tmatrix/output/q\_astrodust\_jori\_P0.20\_Fe0.00\_1.400.dat.gz}, which adds
the 4th block and so is the one that makes $C_{\rm bir,ext}$ nonzero. Either way
the wavelength coverage is the whole axis. Do not confuse this with the five-band
aligned scattering table of \S\ref{sec:aligned} --- the polarized
\emph{extinction} optics are on the full grid.

When a model grid is extended shortward of $0.0912\um$ (through \code{lam\_min}),
\code{build\_Cpol} looks for the companion table
\code{q\_astrodust\_jori\_euv\_P0.20\_Fe0.00\_1.400.dat.gz}. \textbf{That table is
not shipped.} The default state is therefore that the extreme-ultraviolet block
stays at the zero it was initialized with, and the reader says so on stderr:
\begin{lstlisting}[style=sh]
build_Cpol: no EUV polarized table (...)
            dichroic extinction and birefringence are zero below 9.120E-02 um
            (zero by omission, not by physics).
\end{lstlisting}

\paragraph{That zero is a documented deficit, not the right value.} It is worth
being blunt about this, because a zero is easy to mistake for an asymptote.
First-principles size integrals
(\code{tmatrix/driver/euv\_polarized\_optics.f90}) put the alignment-weighted
$|C_{\rm pol,ext}|/C_{\rm ext}$ at $1.26\times10^{-3}$ at the $0.0912\um$ seam,
\emph{rising} to $3.64\times10^{-3}$ near 20.6\,eV --- a factor 2.9 --- before
decaying to $1.6\times10^{-4}$ at 100\,eV, with the sign \emph{opposite} to the
optical band; the reversal falls near $0.106\um$, between $0.1072$ and
$0.1059\um$. $|C_{\rm bir,ext}|/C_{\rm ext}$ falls monotonically from
$7.2\times10^{-3}$ at the seam and changes sign near 51\,eV. Zero would be the
correct limit only for $x\gg1$ \emph{and} $|m-1|\ll1$, and $|m-1|$ is $0.765$ at
the seam and below $0.1$ only above ${\sim}50$\,eV while the grains carrying the
signal sit at $x = 5$--$50$. A sphere has \emph{exactly} zero dichroic extinction
and exactly zero birefringence, so the scalar extreme-ultraviolet optics
(\code{q\_astrodust\_mod}, Mie on the volume-equivalent sphere) could not have
supplied these entries either --- the shape has to be kept to get them at all.

\paragraph{When the companion table is present.} \code{build\_Cpol} reads it and
interpolates in $\log\lambda$ and $\log a$, as \Cabs\ and $C_{\rm sca}$ are
interpolated. A grid running outside the wavelengths the table covers
($0.0124$--$0.0912\um$) is rejected with \code{status = 9} rather than filled by
extrapolation. Generate the table with the \code{euv} mode of
\code{run\_q\_jori.x}, or a subset of it with the named-wavelength modes of
\S\ref{sec:namedlam}.

\paragraph{Why the axis stops at $0.0124\um$ (100\,eV).} Above $x = 60$ the only
description left is the opaque geometric-optics limit, whose absorption is the
Fresnel surface integral described in \S\ref{sec:firstprinciples} and which needs
the chord optical depth $4\,\mathrm{Im}(m)\,x$ to be large. For DH21 astrodust
that quantity is $3.7$ at $x = 50$ and 100\,eV (2.6\% transmission), and it falls
below $1$ by 200\,eV, where the grain is becoming X-ray transparent and neither
the extinction paradox nor the opaque Fresnel absorption applies. The floor is
set where the physics of every branch used is still valid.

\subsection{Certifying the T-matrix above $x = 50$}\label{sec:certify}

The extreme ultraviolet pushes size parameters up, so the $x > 50$ cut of
\S\ref{sec:firstprinciples} would throw away most of the band. In the \code{euv},
\code{lam} and \code{lamfile} modes the T-matrix is therefore attempted up to
$x = 60$ (\code{X\_TM\_MAX}), under a two-setting certification: every node with
$x\in(50,60]$ is solved twice, at the production
$(\code{DDELT},\code{NDGS}) = (10^{-3},2)$ and at a coarser-quadrature
$(3\times10^{-3},3)$, and is \emph{kept only} if the two agree to
\code{TOL\_CHK} $= 5\%$ on every channel. A node that fails falls to geometric
optics. Above $x = 60$ geometric optics is used directly. The plain and
\code{range} sweeps keep the older $x > 50$ rule unchanged, so they still
reproduce the shipped table byte for byte.

Three details are worth recording.

\begin{itemize}
\item \textbf{\code{NDGS} $= 4$ is not usable} as the second setting:
      $\code{NGAUSS} = \code{NDGS}\times N_{\max}$ then exceeds the
      \code{NPNG1} $= 300$ storage limit, so the check would fail for want of
      array space rather than for want of convergence. \code{NDGS} $=3$ is the
      largest that leaves the comparison meaningful.
\item \textbf{\code{IERR} $= 0$ is not sufficient at $x\gtrsim55$.} Mishchenko's
      own convergence flag can report success and still return a wrong answer:
      at $\lambda = 0.0912\um$ and $a = 0.9441\um$ ($x = 65.04$) it returns
      \code{IERR} $= 0$ with a value 50\% away from its neighbors. Measured at the
      seam, $x = 51.7$ and $54.7$ pass the two-setting check while $57.9$, $61.4$,
      $65.0$ and $68.9$ fail it. The certification exists because the engine's own
      flag does not carry this information.
\item \textbf{The polarized channels need an absolute floor.} $Q_{\rm pol}$ is a
      difference of order-unity numbers, so a purely relative comparison becomes
      meaningless where it passes through zero. \code{Q\_POL\_FLOOR} $= 10^{-3}$
      sets the absolute scale below which the polarized channels are not required
      to agree; after the size integral that corresponds to
      $|C_{\rm pol}|/C_{\rm ext}\sim5\times10^{-4}$, about 1\% of the $V$-band
      dichroism.
\end{itemize}

\paragraph{What the certification costs.} Leaving the $x > 60$ nodes at the
geometric-optics value, which carries $Q_{\rm pol,ext} = Q_{\rm bir,ext} = 0$
exactly, loses part of the \emph{dichroic extinction}: nothing at the seam,
${\sim}8\%$ of the band total at $0.031\um$, ${\sim}20\%$ at $0.022\um$ and
${\sim}46\%$ at $0.0124\um$. The true value lies between the certified figure
(a lower bound, $1.06\times10^{-4}$) and a $1/x$ continuation (an upper bound,
$3.14\times10^{-4}$). No extrapolation is put in place of the missing part; the
absolute level is already small there, about 1\% of the optical-band signal.
\Cpol\ --- the absorption dichroism, which drives polarized \emph{emission} ---
has no such gap, because the geometric-optics limit obtains it from the
opaque-grain Fresnel surface integral.

%====================================================================
\section{Aligned-grain scattering matrix and the extinction matrix}
\label{sec:aligned}

Section~\ref{sec:firstprinciples} regenerated the orientation-resolved
\emph{cross sections}. This section adds the angle-resolved object they left out:
the fixed-orientation Mueller matrix $Z(\theta_i;\theta_s,\phi)$ of the aligned
astrodust spheroid, and the $4\times4$ extinction matrix $K(\theta_i)$ that goes
with it. Between them they complete the material side of the Peest-formalism
contract --- SEDust returns every quantity in matrix form; the RT code (MoCafe,
whose consumption side is not yet built) does the rotations, sampling, peel-off,
and $\exp(-K\tau)$ transfer. The design is set out in
\code{docs/tmatrix\_aligned\_scattering\_plan.md}; this section records how each
piece was built and why.

\subsection{The engine}\label{sec:aligned:engine}

\code{tmatrix/driver/scattering\_matrix\_oriented.f90} returns the
fixed-orientation Mueller matrix $Z$ of the DH21 oblate spheroid ($b/a=1.4$) for
one incidence direction. Two regimes, one output convention.

\paragraph{T-matrix regime, one \code{AMPL} call.} With the T-matrix already
converged and held in the workspace by
\code{tmatrix\_eval\_scattering\_matrix} --- the storage separation of
\S\ref{sec:firstprinciples} that made the oriented cross sections
possible --- a single call to Mishchenko's \code{AMPL} evaluates the $2\times2$
complex amplitude matrix $\mathbf{S}=[[S_{vv},S_{vh}],[S_{hv},S_{hh}]]$ for the
chosen incidence and scattering angles; the multipole truncation order and the
wavelength come from the returned \code{scatmat} and from the workspace, never
from the caller. The 16 phase-matrix elements then follow
from the standard bilinears in $\mathbf{S}$, copied \emph{verbatim} from
Mishchenko's \code{reference/upstream/ampld.lp.f} (its
\code{Z11}$\ldots$\code{Z44}, Eqs.\
(13)--(29) of Appl.\ Opt.\ \textbf{39}, 1026, 2000) so the phase matrix can be
checked line by line against the source --- which is why that unmodified file is
kept in the tree even though nothing compiles it. Copying rather than rederiving is
deliberate: the bilinears carry sign and conjugation conventions that are easy to
transcribe wrongly, and the whole point of anchor~C below is that this transcription
is exact.

\paragraph{Rayleigh regime, analytic dipole.} Below $x=0.1$ the same $Z$ is the
electric-dipole matrix, $S_{pq}=k^2\,(\hat e_p^{\rm sca}\!\cdot\!\boldsymbol\alpha
\!\cdot\!\hat e_q^{\rm inc})$ with the spheroid polarizability
$\boldsymbol\alpha=\mathrm{diag}(\alpha_b,\alpha_b,\alpha_a)$, then the same
bilinears. The polarizability comes from \code{spheroid\_dipole\_polarizability},
a routine extracted from \code{rayleigh\_limit} (\S\ref{sec:firstprinciples})
without changing an operation, so the dipole cross sections that
\S\ref{sec:firstprinciples} already validated and the dipole matrix here rest on
one implementation of $\alpha_a,\alpha_b$.

\paragraph{The Stokes basis, stated once.} Both regimes use Mishchenko's
meridional basis of each propagation direction \emph{in the grain frame}, whose
$z$ axis is the alignment axis: $\hat e_1=\hat\theta$ (call it $v$), $\hat
e_2=\hat\phi$ ($h$), with $Q=I_v-I_h$. This is exactly the basis \code{AMPL}
rotates its amplitudes into, so the analytic dipole and the T-matrix matrix share
one convention with no adaptor between them. At $\theta_i=90^\circ$ the
$v$-polarization is $\mathrm{jori}=2$ ($\mathbf{E}\parallel\hat a$) and $h$ is
$\mathrm{jori}=3$ ($\mathbf{E}\perp\hat a$), which is what ties the extinction
matrix back to the jori cross sections of \S\ref{sec:source}.

\subsection{Symmetries, verified not assumed}\label{sec:aligned:sym}

Three exact symmetries reduce the stored ranges and, more usefully, are numerical
checks on the engine (anchor~E). Each was confirmed on a grid of
$(\theta_i,\theta_s,\phi)$ before being relied on.

\begin{itemize}
\item \textbf{Azimuth mirror.} The plane through the axis and the incidence
      direction is a symmetry plane, so $\phi\to360^\circ-\phi$ leaves $Z$
      unchanged except that the two off-diagonal $2\times2$ blocks (elements
      $13,14,23,24,31,32,41,42$) flip sign. Store $\phi\in[0,180^\circ]$; residual
      $3.5\times10^{-7}$.
\item \textbf{Equatorial mapping.} The oblate spheroid is symmetric under
      $z\to-z$, so $\theta_i\to180^\circ-\theta_i$ maps onto the stored range with
      $\theta_s\to180^\circ-\theta_s$, $\phi$ \emph{unchanged}, and the same
      off-diagonal-block sign flip,
      \[
        Z(180^\circ-\theta_i;\,180^\circ-\theta_s,\,\phi)\;=\;S\odot
        Z(\theta_i;\,\theta_s,\,\phi),\qquad
        S=\begin{pmatrix}+&+&-&-\\ +&+&-&-\\ -&-&+&+\\ -&-&+&+\end{pmatrix},
      \]
      the diagonal blocks keeping their sign. Store $\theta_i\in[0,90^\circ]$;
      residual $8.8\times10^{-16}$. Physically this also encodes that the
      alignment axis is headless.
\item \textbf{Axisymmetry at $\theta_i=0$.} There $Z$ must be independent of
      $\phi$ up to a frame rotation, $Z(0;\theta_s,\phi)=Z(0;\theta_s,0)\,
      R(\phi)$, the $\phi=0$ matrix being six-element block-diagonal. This is a
      check, not a storage saving.
\end{itemize}

The library reconstructs the unstored octants from these relations, applying the
$\pm1$ sign once to the interpolated matrix rather than to eight interpolation
corners (\S\ref{sec:aligned:lib}); because each reflection contributes a constant
that is the same for every corner of a cell, it factors out of the linear
interpolation exactly.

\subsection{The generator: one code path}\label{sec:aligned:gen}

The shared physics lives in \code{tmatrix/driver/aligned\_population\_optics.f90},
and both the production driver \code{run\_scatmat\_aligned.f90} and the
verification program \code{compare\_scatmat\_aligned.f90} call it. This is on
purpose: the numbers an anchor checks are produced by the exact code the
production file is written from, so a passing anchor certifies the shipped path
and not a parallel one.

\paragraph{One solve, three products.} For each (band, size) in the T-matrix
regime a single \code{tmatrix\_eval\_scattering\_matrix} solve serves everything:
the expansion coefficients it returns in the \code{scatmat} feed the
random-orientation accumulation (the $F$ block below), the T-matrix it leaves in
the workspace feeds the \code{AMPL} sweep over the
$Z$ nodes, and its forward amplitudes give the extinction-matrix elements. No
quantity is recomputed.

\paragraph{Regime split and the $x>50$ guard.} The size-parameter regimes match
\code{run\_scatmat.f90}: $x<0.1$ analytic dipole, $0.1\le x\le50$ T-matrix. Above
$x=50$ no converged oriented matrix exists, so the bin is \emph{omitted from the
aligned products} and instead enters $F_{\rm tot}$ as its random-orientation
geometric-optics matrix, i.e.\ scattered as unaligned. Its $C_{\rm sca}$ is
tallied, and the driver stops hard if the omitted fraction of the scattering
weight exceeds $10^{-6}$. It never does for the shipped bands: the measured
omitted fraction is at most $3.1\times10^{-15}$ (at $0.36\um$) and exactly zero at
$0.64$ and $0.79\um$. The omission is thus provably negligible rather than
assumed so.

\paragraph{The threading contract: one active caller for each workspace.} The
only parallel region is the loop over $Z$ nodes inside one size. \code{AMPL} and
\code{VIGAMPL} take the workspace \code{intent(in)}: they read the stored T-matrix
and otherwise write local automatic arrays, so distinct nodes are independent and
each grid slot is written once --- which is why a 1-thread and an $N$-thread run
are bitwise identical. What this rests on is no longer a compile flag but the
shape of the library. There is no \code{COMMON}, no \code{EQUIVALENCE} and no
mutable \code{SAVE} anywhere in it, and every working array lives in the
caller-owned \code{tmatrix\_workspace\_t}, so the contract fits in one line:
\emph{one active caller for each workspace, and distinct workspaces may be used
concurrently}. Reading is the exemption this loop depends on --- an amplitude
evaluation only reads, so several threads may evaluate against a single workspace
at once, and that is exactly what happens here: the T-matrix is solved once,
serially, and only the scattering directions are swept in parallel. Giving each
thread a workspace of its own would not be the safer choice but the wrong one, a
fresh workspace holding no T-matrix and the evaluation returning \code{IERR = 2}.

The whole library is compiled \code{-frecursive}, so large locals stay on the
caller's stack. One file, \code{tmatrix\_core.f90}, additionally carries
\code{-fno-inline}, and that flag has nothing to do with threads: inlining moves
where the compiler contracts multiply-add pairs, which moves the last bits of the
T-matrix convergence test, and without it 2 of the 6258 rows of the shipped
astrodust table shift by ${\sim}10^{-6}$. Lifting the largest scratch arrays out
of the call chain and into the workspace also brought the stack requirement of the
deepest path down from about 3.1\,MB to about 190\,kB; what now scales with the
thread count is the workspace itself, 83\,926\,096 bytes (80.04\,MiB) each.

The lesson worth carrying forward is the reverse of the one this paragraph used to
record. Deciding thread-safety file by file through compile flags holds only as
long as someone remembers which file is which. Removing the shared state and
handing the storage to the caller makes the contract a property of the source
instead, and one sentence long.

\subsection{Products and the $\eta$ contract}\label{sec:aligned:products}

Per band the file (\code{output/scatmat\_aligned\_astrodust\_P0.20\_Fe0.00\_1.400.dat})
carries three blocks, all per H.

\begin{itemize}
\item \textbf{$K$ block}, on the $\theta_i$ grid: $C_{\rm ext,al}$, the dichroic
      $C_{\rm pol,al}$, and the birefringent $C_{\rm bir,al}$ from the forward
      amplitudes, plus $C_{\rm sca,al}$ as the grid closure $\int Z_{11}\,
      d\Omega$. Because these come from the forward amplitude at each $\theta_i$,
      the extinction matrix is \emph{exact at every propagation angle}. This
      replaces the $\sin^2\psi$ interpolation between the $\psi=0$ and $90^\circ$
      values that the three-orientation jori table forces, and reduces to that law
      analytically in the Rayleigh limit (anchor~G).
\item \textbf{$F$ block}: $F_{\rm tot}$ (every grain, identical to
      \code{run\_scatmat.f90}) and $F_{\rm ref}$ ($f_{\rm align}$-weighted), the
      six-element random-orientation matrices, $\alpha_1$-normalized so
      $\tfrac12\int F_{11}\,d\cos\Theta=1$ (equivalently $\int F_{11}\,d\Omega=
      4\pi$); the header gives $C_{\rm sca,tot}$, $C_{\rm sca,ref}$,
      $C_{\rm ext,tot}$, $C_{\rm ext,ref}$, from which the absolute differential
      matrix is $C_{\rm sca}F/(4\pi)$ ($F_{\rm tot}$ with $C_{\rm sca,tot}$,
      $F_{\rm ref}$ with $C_{\rm sca,ref}$).
\item \textbf{$Z$ block}: $Z_{\rm al}=\int da\,n(a)\,f(a)\,Z(a)$, 16 elements,
      $\um^2\,\mathrm{sr}^{-1}$ per H, on
      $\theta_i=0(5)90^\circ$ (19) $\times$ $\theta_s=0(1)180^\circ$ (181)
      $\times$ $\phi=0(5)180^\circ$ (37).
\end{itemize}

\paragraph{The $\eta$ contract.} The file header documents, and the library
enforces, the exact rule for varying the alignment degree cell to cell by a single
scalar $\eta$ (the local alignment scale, $f_{\rm cell}=\eta f_{\rm ref}$). The
aligned matrix scales linearly, $Z_{\rm al,cell}=\eta Z_{\rm al}$; the extinction
matrix likewise, $K_{\rm al,cell}=\eta K_{\rm al}$; the unaligned remainder
scattering matrix is $Z_{\rm unal}=[C_{\rm sca,tot}F_{\rm tot}-\eta C_{\rm sca,ref}
F_{\rm ref}]/(4\pi)$ in absolute units, whose $\int Z_{11}\,d\Omega=C_{\rm sca,tot}
-\eta C_{\rm sca,ref}$; and the unaligned
population adds an isotropic $C_{\rm ext,tot}-\eta C_{\rm ext,ref}$ to the
extinction (its polarized and birefringent parts vanish). All of this is exact by
the linearity of the size integral in $f_{\rm align}$, so three stored integrals
give the scattering optics of any $\eta$ with no re-integration.

\subsection{The library: two strictly separated layers}\label{sec:aligned:lib}

\code{sed/src/scatmat\_aligned.f90} (module \code{scatmat\_aligned\_mod}) serves
the table to a polarized RT host with the same lifecycle as the rest of the
library: initialize once, query along the path.

\paragraph{Initialization (serial, once).} \code{load\_scatmat\_aligned(path,
status)} parses the table into module storage and \code{free\_scatmat\_aligned}
releases it. The resident cost is dominated by \code{scm\_Z}, whose size is
$n_{\theta_i}n_{\theta_s}n_\phi\times16\times n_{\rm band}$ doubles: for the
production grid ($19\times181\times37$, five bands) that is about 81\,MB, linear
in the number of bands, with the $F$ and $K$ arrays adding under 0.1\,MB. Queries
are trilinear interpolations --- a few hundred floating-point operations per
scattering event, with no wavelength search once \code{scatmat\_band} has mapped
the wavelength to a band. It is wired into
\code{sed\_init} / \code{build\_astrodust} through an optional \code{scatmat\_path}
argument --- exactly the \code{qpol\_path} pattern of \S\ref{sec:firstprinciples},
except that a \code{scatmat\_path} that fails to load is an error (\code{status =
3}), because a host that asked for the table depends on it, whereas a missing
polarized-optics table only disables polarized emission. This layer is not
thread-safe and is called from serial code.

\paragraph{Path queries (pure reads, OpenMP-safe).} Six routines write only their
own outputs and read the loaded arrays, so they are safe to call concurrently from
photon threads:
\begin{itemize}
\item \code{scatmat\_band(lambda\_um, iband, exact)} --- nearest stored band, so
      the hot path takes an integer and no wavelength search runs per event;
\item \code{extinction\_matrix\_aligned(iband, theta\_i, eta, kmat)} --- the
      $4\times4$ matrix with diagonal $\eta\,C_{\rm ext,al}$,
      $K(1,2)=K(2,1)=\eta\,C_{\rm pol,al}$, $K(3,4)=\eta\,C_{\rm bir,al}$,
      $K(4,3)=-\eta\,C_{\rm bir,al}$, $\theta_i$ folded into $[0,90^\circ]$ because
      $K$ is even under the equatorial reflection;
\item \code{mueller\_matrix\_aligned(iband, theta\_i, theta\_s, phi, z)} ---
      trilinear interpolation plus the symmetry reconstruction of
      \S\ref{sec:aligned:sym}, at the reference alignment ($\eta=1$; the caller
      scales);
\item \code{mueller\_matrix\_random(iband, big\_theta, f\_tot, f\_ref)} --- the two
      six-element random-orientation matrices;
\item \code{mueller\_matrix\_total(iband, theta\_i, theta\_s, phi, eta, z)} --- the
      recommended host call: the \emph{absolute} combined matrix
      $z=\eta Z_{\rm al}+L(\pi-\sigma_2)F_{\rm unal}(\Theta)L(-\sigma_1)/(4\pi)$
      with $F_{\rm unal}=C_{\rm sca,tot}F_{\rm tot}-\eta C_{\rm sca,ref}F_{\rm ref}$,
      $\cos\Theta=\cos\theta_i\cos\theta_s+\sin\theta_i\sin\theta_s\cos\phi$ and the
      meridional rotation angles $\sigma_1,\sigma_2$ from
      \code{meridional\_scattering\_angles} (closed \code{atan2} forms, poles
      included), so a host need not reassemble the mixture from the arrays;
\item \code{scattering\_cross\_sections(iband, theta\_i, eta, csca\_aligned,
      csca\_unaligned [, csca\_pol\_aligned])} with $\code{csca\_unaligned}=
      C_{\rm sca,tot}-\eta\,C_{\rm sca,ref}$ and the optional
      $\code{csca\_pol\_aligned}=\eta\,C_{\rm sca,pol,al}(\theta_i)=\eta\int
      Z_{12}\,d\Omega$, the aligned matrix's polarized scattering cross section
      (Peest et al.\ 2023; the random remainder integrates to zero), on the same
      $[0,90^\circ]$ fold.
\end{itemize}
\code{scm\_csca\_pol\_al(nti,nband)}, the polarized scattering cross section
$\int Z_{12}\,d\Omega$ that \code{csca\_pol\_aligned} returns, is formed at load
time with the generator's own closure quadrature --- the same rule that closes
$\int Z_{11}\,d\Omega$ to $C_{\rm sca,al}$ --- so it needs no extra file column.
The loaded grids and arrays are also exposed \code{public, protected}, so a host
that prefers to inline its own lookups can take them once at startup and never
call back. The whole cell dependence enters through just two scalars the host
already holds: $\eta$, and $\theta_i=\arccos(\hat k\!\cdot\!\hat B)$ from one dot
product.

\paragraph{The alignment guard (status 4).} The $K$ and $Z$ arrays were integrated
over size once, under the profile recorded in the header; the sanctioned runtime
variation is $\eta$, not a change of profile. So when a loaded scattering table is
present and the host calls \code{dust\_set\_alignment} or
\code{dust\_set\_alignment\_profile} with a profile that differs from the recorded
one, the setter returns a \emph{new} status code~\code{4}. It is non-fatal: the
$f_{\rm align}$ update still completes (the emission and $C_{\rm pol,ext}$ channels
re-integrate live and honor the new profile), and the code only signals that the
loaded scattering matrices no longer correspond to the model's alignment. A
different profile requires regenerating the table
(\code{run\_scatmat\_aligned.x profile=FILE}), which takes minutes. A tabulated
profile never matches the recorded analytic one and so always raises the flag.

\subsection{Birefringence through \code{dust\_extinction}}\label{sec:aligned:cbir}

\code{dust\_extinction} gained an optional \code{Cbir\_ext} output, the
$f_{\rm align}$-weighted size integral of the jori table's 4th-block
birefringence,
\[
  C_{\rm bir,ext}(\lambda)=\sum_a {\rm d}n_{\rm Ad}(a)\,
     C_{\rm bir,ext}(\lambda,a)\,f_{\rm align}(a),
\]
mirroring \code{Cpol\_ext} exactly. It is zero when the loaded table carries no
4th block (as the release table does not), so the default path is unchanged, and
it is the same birefringence the aligned $K$ block carries at $\theta_i=90^\circ$,
reached here without loading the aligned scattering table. This closes the gap
noted in the previous release, that \code{dust\_extinction} did not yet return the
birefringence.

\subsection{Verification}\label{sec:aligned:verif}

Every anchor of the plan passes (Table~\ref{tab:alignverif}). The engine anchors
(A--H) run at single sizes through \code{compare\_scatmat\_aligned.x} --- anchor~H
adds the orientation-averaged amplitude engine at general geometry
($\theta_i=40^\circ$, $\theta_s=70^\circ$, $\phi\in\{30,120,250\}^\circ$) against
$L(\pi-\sigma_2)F(\Theta)L(-\sigma_1)$, which a wrong $\sigma$ sign fails at order
unity to order 30. The library checks, including the failing-tolerance
certification tests [3a]--[3e] of \code{test\_scatmat\_aligned.x} --- the absolute
$4\pi$ closure, the $\int Z_{12}$ polarized cross section, the
\code{mueller\_matrix\_total} closure, and the $\theta_i=0$ rotation sign --- run
on the production table.

\begin{table}[h]
\centering
\small
\begin{tabular}{@{}llp{0.30\textwidth}l@{}}
\toprule
& Check & What it certifies & Result \\
\midrule
\multicolumn{4}{@{}l}{\emph{Engine (single size)}}\\
A & closure $\int Z_{11}\,d\Omega$ vs $C_{\rm sca}(\mathrm{jori})$ & phase-matrix normalization & $2.7\times10^{-7}$ \\
B & optical theorem vs $C_{\rm ext}(\mathrm{jori})$ & forward-amplitude geometry & $2\times10^{-14}$ \\
C & orientation average vs GSP $F(\Theta)$ & the bilinears, incl.\ $F_{12},F_{34}$ signs & $1.4\times10^{-6}$ \\
D & dipole vs closed form & Rayleigh implementation & $2\times10^{-16}$ \\
  & dipole vs T-matrix at $x\!\sim\!0.1$ & regime continuity, $|\delta Z|/Z_{11}$ & $0.5$--$1.1\%$ \\
E & $\phi$ mirror & azimuth symmetry & $3.5\times10^{-7}$ \\
  & equatorial mapping & $z\to-z$ symmetry & $8.8\times10^{-16}$ \\
G & Rayleigh $K(\theta_i)$ vs $\sin^2$ law & extinction-matrix angle law (exact) & $3.3\times10^{-16}$ \\
  & T-matrix $K(\theta_i)$ vs $\sin^2$ law & at $x=0.137$ & $6.3\times10^{-4}$ \\
H & general geometry & amplitude engine vs $L(\pi\!-\!\sigma_2)F(\Theta)L(-\sigma_1)$, $\theta_i{=}40^\circ,\theta_s{=}70^\circ$; wrong $\sigma$ sign fails O(1)--O(30) & $4.1\times10^{-5}$ \\
\midrule
\multicolumn{4}{@{}l}{\emph{Generator (size-integrated)}}\\
F & $K$ vs independent jori integrals & $C_{\rm ext}$ & $1.6\times10^{-6}$ \\
  &                                   & $C_{\rm pol}$ & $2.0\times10^{-5}$ \\
  &                                   & $C_{\rm bir}$ & $1.3\times10^{-7}$ \\
\midrule
\multicolumn{4}{@{}l}{\emph{Library (production table)}}\\
  & reader round-trip & parse vs re-emit & bitwise 0 \\
  & symmetry reconstruction & unstored octants & bitwise 0 \\
  & $\eta$ algebra & linear scaling & residual 0 \\
  & $K(90^\circ)$ vs jori vs \code{dust\_extinction} & $C_{\rm pol}$ / $C_{\rm bir}$ & $1.6\times10^{-5}$ / $8.0\times10^{-5}$ \\
  & \code{rt\_example} closures & end-to-end consumption & $0.997$--$1.001$ \\
{[3a]} & random absolute closure & $(4\pi)^{-1}\!\int C_{\rm sca,tot}F_{11}\,d\Omega=C_{\rm sca,tot}$ & $3.5\times10^{-4}$ \\
{[3b]} & stored $Z_{11}$ re-integration & vs the table's $C_{\rm sca,al}$ column (ASCII rounding) & $3.3\times10^{-6}$ \\
{[3c]} & $C_{\rm sca,pol}$ vs $\int Z_{12}\,d\Omega$ & routine vs direct; $|C_{\rm sca,pol}|/C_{\rm sca,al}\le0.101$ & exact \\
{[3d]} & \code{mueller\_matrix\_total} closure & $\int z_{11}=\eta C_{\rm sca,al}+C_{\rm sca,tot}-\eta C_{\rm sca,ref}$; $\int z_{12}=\eta C_{\rm sca,pol}$; $\eta\in\{0,0.37,1\}$ & $2.6\!\times\!10^{-4}$ / $7\!\times\!10^{-7}$ \\
{[3e]} & rotation sign at $\theta_i=0$ & $Z(0;\theta_s,\phi)=Z(0;\theta_s,0)\,L(-\phi)$ & $6.8\times10^{-5}$ \\
\bottomrule
\end{tabular}
\caption{Verification of the aligned scattering matrix, its extinction matrix, and
the library API. Engine and generator anchors are those of the plan's \S5, plus
anchor~H (the amplitude engine at general geometry) and the failing-tolerance
library certification tests [3a]--[3e] of \code{test\_scatmat\_aligned.x}; the
library checks run on the shipped production table.}
\label{tab:alignverif}
\end{table}

\paragraph{Cost of the production run.} The five UBVRI bands (0.36, 0.44, 0.55,
0.64, 0.79\um) take about 8 minutes at 32 threads (about 26 minutes at 8); the
file is 138\,MB, 38\,MB gzipped. The existing regressions are unchanged:
\code{calc\_polext} reproduces the release to a median of $0.0296\%$ as before,
and every \code{sed}, \code{mc} and \code{tmatrix} target builds clean.

\subsection{Which wavelengths the table carries}\label{sec:aligned:bands}

The shipped table holds five bands, $0.36$, $0.44$, $0.55$, $0.64$ and
$0.79\um$ (\code{BANDS} in \code{run\_scatmat\_aligned.f90}), on the grid
$\theta_i = 0(5)90^\circ$ (19) $\times$ $\theta_s = 0(1)180^\circ$ (181) $\times$
$\phi = 0(5)180^\circ$ (37). That is the intended scope: polarized transfer is
essentially never run on a whole wavelength grid, and five optical bands cover
the update this work is for. What matters is that a wavelength one later needs
can be computed, not that everything is precomputed now.

\code{run\_scatmat\_aligned.x} does already take wavelengths on its command line
(\code{./run\_scatmat\_aligned.x 0.44 0.55 0.79}), so a different set is one run
away. Three limitations of that path should be stated plainly rather than
discovered:
\begin{itemize}
\item \textbf{The output stem is fixed} (\code{f\_stem//'.dat'}), so a custom run
      \emph{overwrites} the shipped five-band table. Move or rename the shipped
      file first.
\item \textbf{There is no merge mode.} A band cannot be added to an existing
      table; every requested band is recomputed from scratch into a new file
      (the five together are the 8-minute run above).
\item \textbf{There is no size-window split and no large-$x$ certification.} The
      driver has no \code{ja=} equivalent of \S\ref{sec:namedlam}, and it applies
      only the plain \code{X\_LARGE} $= 50$ cut, without the two-setting
      cross-check of \S\ref{sec:certify}.
\end{itemize}
None of these is hard to lift; none of them is done, and the wavelengths in hand
have not required it.

\subsection{Why the library reads tables and does not solve}\label{sec:notmatrix}

The library never runs a T-matrix solve at run time. The reason is not the size of
the code --- the engine is already in the tree and the drivers link it --- but
that convergence cannot be certified at $x\gtrsim55$ from inside a transport loop:
as \S\ref{sec:certify} records, Mishchenko's \code{IERR} can report success on a
node that is 50\% wrong, and the check that catches it costs a second solve at
different settings. A silently wrong optic is worse than a missing one, and the
transport loop has nowhere to report a failure to. The tables are generated
offline, where a node can be rejected and the rejection counted in the file
header. This is the current policy rather than a permanent decision: the T-matrix
core itself is expected to be revised, and a certification cheap enough to run
inline would change the trade.

%====================================================================
\section{Status codes, and the relation to v1.00}\label{sec:status}

\subsection{\code{sed\_init} / \code{build\_astrodust} status codes}

Both take an optional \code{status}; $0$ is success, and when the argument is
absent the readers keep their message-and-stop behavior instead. The v1.20 codes
are those of Table~\ref{tab:status}.

\begin{table}[h]
\centering
\small
\begin{tabular}{@{}cp{0.78\textwidth}@{}}
\toprule
\code{status} & Meaning \\
\midrule
1 & $Q$-table load failed \\
2 & size-distribution load failed \\
3 & aligned scattering table load failed (only when \code{scatmat\_path} was
    supplied --- a host that asked for it depends on it) \\
4 & a polarized $Q$ table requested explicitly through \code{qpol\_path} could
    not be read (an implicit default degrades gracefully instead) \\
5 & \code{load\_polarized\_optics = .false.} combined with an explicit
    polarized-optics path --- a contradictory request \\
6 & astrodust dielectric function load failed (extreme-ultraviolet band only) \\
7 & \code{lam\_min} below the astrodust dielectric function's own shortest
    wavelength (extreme-ultraviolet band only) \\
8 & the extreme-ultraviolet polarized table could not be read while the grid
    reaches into that band and polarized optics were asked for \\
9 & the extreme-ultraviolet band of the grid runs outside the wavelengths the
    companion polarized table covers ($0.0124$--$0.0912\um$) \\
\bottomrule
\end{tabular}
\caption{Status codes of \code{sed\_init} and \code{build\_astrodust} in v1.20.}
\label{tab:status}
\end{table}

Code~8 is declared for an unreadable extreme-ultraviolet polarized table. Note
what the current code actually does with an \emph{absent} companion table: that is
not an error --- \code{build\_Cpol} reports it on stderr and leaves the band zero
(\S\ref{sec:euvpol}) --- so the status a caller sees in practice is $9$, raised
when a table that \emph{was} read does not cover the requested grid.

\code{dust\_set\_alignment} and \code{dust\_set\_alignment\_profile} have their own
independent numbering: $1$, $2$, $3$ for an out-of-range \code{a\_align},
\code{alpha\_align} or \code{f\_max}, and $4$ for the non-fatal alignment guard of
\S\ref{sec:aligned:lib}. Do not read a $4$ from a setter as the \code{qpol\_path}
failure above.

\subsection{v1.20 is v1.00 plus polarization}

The two maintained branches are a superset relation, not a fork. Everything scalar
is identical: the same builders (\code{build\_astrodust}, \code{build\_dl07},
\code{build\_zubko}, \code{build\_from\_files}), the same \code{dust\_emission} and
\code{dust\_extinction}, the same \code{lam\_min} extreme-ultraviolet extension,
and the same \code{calc\_kext.x} products. What v1.20 adds is the polarized
material: \code{Cpol}, \code{Cpol\_ext}, \code{Cbir\_ext}, \falign, the aligned
scattering table and its library layer. The only incidental difference is the
status numbering --- what v1.00 calls $3$ and $4$ (the dielectric function load
and the \code{lam\_min} floor) are $6$ and $7$ here, because $3$, $4$ and $5$ were
taken by the polarized paths. A scalar-only v1.20 run
(\code{load\_polarized\_optics = .false.}) returns the scalar \code{Cext},
\code{Cabs}, \code{Csca}, \code{gbar} and the total SED bit-identical to the
polarized build at zero alignment.

%====================================================================
\section{Limits and extension paths}

These are limits of the published optics, not of the code.

\paragraph{Birefringence is now computed; the release table has none.} The
release table stores $Q_{\rm ext}$, $Q_{\rm abs}$ and $Q_{\rm sca}$ only, no phase
retardation, so from it $C_{\rm circ}=0$. The first-principles route of
\S\ref{sec:firstprinciples} already evaluates the fixed-orientation amplitude
matrix (Mishchenko's \code{AMPL}, in
\code{tmatrix/src/fixed\_orientation\_amplitude.f90}) and used
only its imaginary forward part, for extinction. Keeping the real part as well
costs nothing --- it comes from the same amplitude call --- and gives the
birefringence, the phase retardation that couples $U$ and $V$. It is stored as an
optional 4th block of the regenerated table, $Q_{\rm re}(\mathrm{jori})=
(4\pi/k)\,\mathrm{Re}[S_{\rm fwd}]/(\pi a^2)$, the real twin of $Q_{\rm ext}$, from
which $C_{\rm bir}=\tfrac12(Q_{{\rm re},3}-Q_{{\rm re},2})\,\pi a^2$. The reader
exposes it as \code{qbir\_ext} with a \code{has\_bir} flag that stays \code{.false.}
for an older 3-block table, so nothing about the default path changes. No astrodust
reference for the birefringence exists, so it is certified internally by
Kramers-Kronig: $\mathrm{Re}$ and $\mathrm{Im}$ of the forward-amplitude difference
are a Hilbert pair, and the directly computed birefringence agrees with the Hilbert
transform of the already-validated dichroism to a median $\sim0.1\%$ over the
interior of the grid (the transform itself self-tests to $2\times10^{-5}$ on a
Lorentz oscillator; the residual is the finite-range truncation). What remains is to
return it through \code{dust\_extinction} and to consume it in the $U\leftrightarrow V$
transfer term, both on the host side.

\paragraph{The aligned scattering matrix is now computed.} The release gives the
integrated $Q_{\rm sca}$ with no angular information, so scattering by an aligned
spheroid could not be modeled from it. It is now modeled from first principles:
the fixed-orientation Mueller matrix $Z(\theta_i;\theta_s,\phi)$ of the aligned
astrodust spheroid is built from the same fixed-orientation amplitude
(\code{AMPL}) used above for extinction, size-integrated over the
astrodust distribution with the alignment weight, and paired with the
$4\times4$ extinction matrix $K(\theta_i)$ from the forward amplitudes. It is
written for five optical bands by \code{run\_scatmat\_aligned.x} and served to a
polarized RT host through a two-layer library API. \S\ref{sec:aligned} is the
full account, and \S\ref{sec:aligned:bands} covers the wavelength coverage and
what computing a different band involves.

\paragraph{The random-orientation scattering matrix does exist, in five optical
bands.} This is a different object from the previous paragraph and should not be
confused with it. Dropping the alignment, a randomly oriented spheroid with a
plane of symmetry has six independent scattering-matrix elements rather than
sixteen, and those \emph{are} computed in the tree:
\code{tmatrix/driver/run\_scatmat.f90} accumulates the single-size expansion
coefficients over the HD23 size distribution with $C_{\rm sca}\,{\rm
d}n/n_{\rm H}$ weighting, expands them onto the scattering-angle grid with
\code{SCATMAT\_FROM\_MOMENTS} (Mishchenko's \code{MATR}, in module
\code{scattering\_matrix\_expansion}), and writes
\code{tmatrix/output/scatmat\_astrodust\_P0.20\_Fe0.00\_1.400.dat}: one block per
wavelength, 181 rows at $\theta = 0, 1, \ldots, 180^\circ$, with columns
$F_{11}, F_{22}, F_{33}, F_{44}, F_{12}, F_{34}$. It answers what an unaligned
grain population scatters; it does not answer the aligned question above.

The stored wavelengths are $0.36$, $0.44$, $0.55$, $0.64$ and $0.79\um$, roughly
\textit{UBVRI} --- the same five bands the aligned table carries, and for the same
reason (\S\ref{sec:aligned:bands}). Five bands are the deliverable, not a sweep
left half-finished: polarized transfer is run at a few wavelengths, the $V$ band
first with one or two further optical bands possible, and the whole 1129-point
grid would cost about 1.7 hours spent almost entirely on wavelengths nothing is
waiting for. What the design owes is a way to compute a wavelength when one is
wanted, and that costs a run rather than a code change:
\code{./run\_scatmat.x all} does the full grid, or a list of wavelengths does a
subset, at about $5.5\,$s per wavelength.

\paragraph{Far infrared needs neither.} In that regime the albedo is
essentially zero, so the transfer requires only the extinction matrix and the
emission vector, both fully determined by what is exported. The HAWC+ 154\um\
comparison lies inside this regime, so the far-infrared milestone is reachable
with the present data and should be attempted before any scattering work.

\paragraph{Astrodust only.} DL07 and Zubko have no polarized optics.
\code{build\_dl07} releases any polarized arrays left by a previous astrodust
initialization, so \code{lamI\_pol} returns zeros rather than stale values.

\subsection{Cost of \code{build\_Cpol}}

Loading the orientation-resolved table is the most expensive part of model
construction (Table~\ref{tab:cost}).

\begin{table}[h]
\centering
\begin{tabular}{@{}ll@{}}
\toprule
Item & Size \\
\midrule
compressed file          & 4.1\,MB \\
scratch file when expanded & 17.2\,MB \\
raw arrays \code{qext\_j}, \code{qabs\_j}, \code{qsca\_j} & 13.7\,MB, transient \\
derived \code{qpol\_*}, \code{qran\_*} & 7.6\,MB, kept \\
resident after \code{sed\_init} & 7.6\,MB \\
time                     & about 1\,s of the 3.1\,s model build \\
\bottomrule
\end{tabular}
\caption{Cost of loading the orientation-resolved table. The
orientation-resolved arrays exist only between the read and the formation
of the derived combinations.}
\label{tab:cost}
\end{table}

The orientation-resolved arrays \code{qext\_j}, \code{qabs\_j} and
\code{qsca\_j} are deallocated as soon as the combinations of
Eqs.~(\ref{eq:qpol})--(\ref{eq:qran}) have been formed. They are intermediate
quantities: nothing downstream of \code{build\_Cpol} reads them, and every
later step uses the five derived arrays only. Releasing them removes
13.7\,MB, about two-thirds of what the table would otherwise hold, and leaves
7.6\,MB resident for the rest of the run. The three arrays are \code{private}
to \code{q\_table\_jori}, so no consumer can depend on their surviving the
build.

Measured on \code{main\_astrodust.x}, the maximum resident set size falls from
51232\,kB to 40088\,kB, a reduction of about 10.9\,MB. The peak does not fall by the full 13.7\,MB because it
is a high-water mark: once the orientation-resolved arrays are gone, the
later SED solution stages set the peak instead. The steady-state occupancy
just after \code{sed\_init} does drop by the full 13.7\,MB.

%====================================================================
\section{Reproducing the numbers}

From \code{SEDust/sed/}:
\begin{lstlisting}[style=sh]
make calc_polext.x           # standalone polarized-extinction check
./calc_polext.x              # -> output/polarized_extinction_ours.dat
                             #    lambda, ours, reference, ratio

make calc_kext.x             # full-pipeline extinction table, all dust models
./calc_kext.x astrodust      # -> ../data/kext_astrodust_MW.dat, 8 columns
                             #    ./calc_kext.x astrodust euv       (EUV grid)
                             #    ./calc_kext.x astrodust 0.03      (lam_min [um])
                             #    ./calc_kext.x dl07 | zubko | from_files ...

make libsedust.a             # library, for the emission path
\end{lstlisting}

\code{calc\_polext.x} prints the median and maximum deviation against the
release directly, restricted to $\lambda \ge 0.11\um$. Column 8 of the
extinction table should reproduce it to the ASCII precision quoted in
Table~\ref{tab:verif}; if the two disagree by more than that, the $\log a$
interpolation in \code{build\_Cpol} or the size grid it interpolates onto is
the place to look, since that is the only step separating the two routes.

For the emission fractions of Table~\ref{tab:polfrac}, build a model with
\code{build\_astrodust}, evaluate \code{J\_Mathis} at $U = 1.585$, and call
\code{dust\_emission} with \code{lamI\_pol} requested; the example in \S5.4 of
the user manual is the complete recipe.

%====================================================================
\section{References}

\begin{itemize}
\item Draine, B. T., \& Hensley, B. S. 2021b, ApJ, 919, 65 (DH21b).
      \code{2021ApJ...919...65D}
\item Hensley, B. S., \& Draine, B. T. 2023, ApJ, 948, 55 (HD23).
      \code{2023ApJ...948...55H}
\item Peest, C., Camps, P., Stalevski, M., Baes, M., \& Siebenmorgen, R.
      2017, A\&A, 601, A92.
      \code{2017A\&A...601A..92P}
\item Peest, C., et al.\ 2023, A\&A, 673, A112.
      \code{2023A\&A...673A.112P}
\item Planck Collaboration 2020, A\&A, 641, A12.
      \code{2020A\&A...641A..12P}
\item Reissl, S., Wolf, S., \& Brauer, R. 2016, A\&A, 593, A87 (POLARIS I).
      \code{2016A\&A...593A..87R}
\item Seon, K.-I. 2018, ApJ, 862, 87.
      \code{2018ApJ...862...87S}
\end{itemize}

\end{document}
